
@@M0@@
\uline{(where} @@M1@@\uline{) is a polynomial map. If} @@M2@@\uline{, then} @@M3@@ \uline{has degree} @@M4@@.

\emph{Proof}: Let @@M5@@ be a sequence of elements of @@M6@@ with @@M7@@ (@@M8@@). Then the binomial formula gives, for every @@M9@@:
\sgaStableTarget{sga6.e10.eq.2-1-5}@@M10@@
where the sum on the right is finite because the elements @@M11@@ are nilpotent. But the coefficients @@M12@@ are polynomials with rational coefficients in the variables @@M13@@, of degree @@M14@@, so @@M15@@ is indeed a polynomial map. If @@M16@@, then @@M17@@ for @@M18@@. Thus @@M19@@ has degree @@M20@@, Q.E.D.

\sgaStableTarget{sga6.e10.sec.2.2}\subsection*{2.2.}
Now consider a noetherian scheme @@M21@@ and an element @@M22@@ of @@M23@@ (1.1.1). Associating to every @@M24@@ the element @@M25@@ of @@M26@@ gives a homomorphism of @@M27@@-modules
@@M28@@
which, taking account of (1.3.2), factors through the canonical epimorphism
@@M29@@

